\chapter{Conclusion}
\label{ch:conclusion}

\section{Objectives revisited}

Table~\ref{tab:objectives} maps each objective from \S\ref{sec:objectives} to the
evidence bearing on it. Five were met outright; one was met in a different form
from the one proposed.

\begin{table}[htbp]
\centering
\footnotesize
\caption[Objectives against evidence]{Objectives against the evidence in this
report.}
\label{tab:objectives}
\begin{tabularx}{\textwidth}{|c|>{\raggedright\arraybackslash}p{5.0cm}|X|>{\raggedright\arraybackslash}p{1.7cm}|}
\hline
& \textbf{Objective} & \textbf{Evidence} & \textbf{Status} \\
\hline
O1 & Three agents + summariser & \S\ref{sec:grounding}, Table~\ref{tab:agent-config};
 interests traced to Doc 1 categories & Met \\
\hline
O2 & Deterministic orchestration, validated output & \S\ref{sec:orchestration};
 structural compliance 0.9986 over 2{,}865 responses & Met \\
\hline
O3 & Conditionally verifying court agent & \S\ref{sec:court-design};
 0 of 9 flagged instances Court-originated & Met \\
\hline
O4 & Real documents to structured scenarios & \S\ref{sec:extraction};
 23 cases ingested through the user upload path & Met \\
\hline
O5 & Verified corpus, baselines, ablations & \S\ref{sec:baselines}--\S\ref{sec:ablation};
 7 of 8 leak-free merits dispositions matched & Met \\
\hline
O6 & Something the client can use & \S\ref{sec:risk-screen};
 rule-based screen, not the simulator & Met, reframed \\
\hline
\end{tabularx}
\end{table}

\section{What the research questions turned out to have as answers}

\textbf{RQ1, on prompt engineering without fine-tuning.} Yes broadly, with a
caveat that surfaced only under testing. Agents specified through role prompts and
injected context do sustain distinct, role-appropriate positions: the authority
concedes rarely on opening and about half the time thereafter, the bidder far
less, and each reports the opposite kind of outcome against its own declared
fallback. The caveat is that a prompt developed against one dispute shape does not
generalise to another. Confronted with a case where nobody was ever scored, both
agents invented a scoring dispute, because that was the only script they had.

\textbf{RQ2, on which part of the protocol is load-bearing.} The adjudicator,
and essentially nothing else: removing it deadlocks every run. That is close to
trivial in hindsight, since nothing else in the graph can set the resolution flag,
but it remains the one unambiguous architectural finding here, and it is worth
more than the accuracy figures precisely because those do not distinguish this
system from a single prompt.

\textbf{RQ3, on constraining an LLM adjudicator's arithmetic.} Largely yes, and
the boundary is now mapped rather than assumed. Gating exact computation on the
presence of both a formula and every input it needs, and banning unstated numbers
even when hedged as illustrative, held across 425 qualitative assessments with no
Court-originated fabrication among the nine flagged. It does not hold universally:
the Court invented an illustrative dataset on a case it had misclassified as
numeric, and the arithmetic defect on the client's own scoring table resisted
three prompt revisions. The discipline constrains what the adjudicator computes,
not what premise it computes on.

\textbf{RQ4, on evaluating in the absence of ground truth.} A stack of partial
checks with explicit sample sizes, and a refusal to collapse them into one score.
It needs an architecture check, because a system can agree with reality for
reasons unrelated to its design; fabrication screens, because a correct conclusion
from invented evidence looks identical from outside to an honest one; and a
leakage audit, because a scenario can hand over the answer. It also needs to
report negative results, since the two most informative findings here, the
zero-shot tie and the readability shortfall, both count against the artefact.

\section{Reflection on the approach}

The scope ended up wide: five agent types, a state-machine orchestrator, a
document ingestion path, a streaming API, a web application, a rule-based risk
screen, a 23-case verified legal corpus and ten analysis scripts, held together by
92 tests. The complexity that mattered was none of those individually. It was that
the domain punishes exactly what generative models are best at: a model asked
about a procurement dispute produces something fluent, structured and legally
shaped, and the harder the case, the more confident it sounds. Almost every
decision in Chapter~\ref{ch:methodology} traces back to that: schema-validated
output so malformed responses cannot pass silently, compliance counters so repairs
stay visible, an extraction agent instructed as firmly against dropping real
numbers as against inventing them, a source-length guard because an empty PDF once
produced a fictional dispute, and a verification protocol whose purpose is getting
the adjudicator to say ``I cannot compute this'' when it cannot.

Two things I would do differently. First, build the corpus before the agents.
The synthetic scenarios written early on were AI-authored, which made every result
derived from them circular, and they were deleted along with seven directories of
evaluation output. That cost roughly a month, though not for nothing: a
model-authored scenario cannot falsify a model-based system, and discovering that
early enough to rebuild on real case law is what makes every figure in
Chapter~\ref{ch:evaluation} mean anything. Second, test on documents I had not
chosen much sooner, for the reason \S\ref{sec:reflection} gives.

On execution quality, the structural layer is strong and the measurements say so:
0.9986 clean structured responses across 2{,}865 calls, a green 92-case test
suite, every figure reproducible from committed logs by a named command, and
provenance carried through to the analytics interface so a partial batch cannot be
mistaken for a finished one. The reasoning layer is where the limits sit, and they
are limits this project can state numerically rather than suspect
\cite{dodge2019showyourwork}: a 23.1\% wrong-regime citation rate, a summariser
missing its own accessibility claim by 35 points of Flesch Reading Ease, an outcome
vocabulary that cannot express three real remedies in its own corpus, six scenarios
that hand the Court its answer, and one arithmetic defect prompt engineering never
closed. Each is quantified, traced to a named case or command, and paired with its
fix in \S\ref{sec:future}. Failure modes bounded this precisely are a better
starting point than failures nobody looked for.

\section{Future work}
\label{sec:future}

Five directions follow from findings reported above rather than from speculation.

\textbf{Ground the citations.} A 23.1\% wrong-regime rate is the strongest
argument anywhere in this report for the retrieval layer that was scoped out
earlier. Retrieval \cite{lewis2020rag} over the Act and the TCC guidance targets
the failure observed: the model does not know which regime governs a given year,
and no instruction can hand it that knowledge.

\textbf{Check premises the way arithmetic is checked.} The Court agent verifies
numbers against the scenario record but accepts factual assertions from the
negotiating agents without checking them against that same record, and the
Faraday result shows what that costs. A grounding check on negotiating-agent
statements, built on the machinery the numeric screen already provides, is a
concrete next step rather than an open-ended one.

\textbf{Strip the disposition at ingestion.} Six of fourteen merits scenarios
state what the real court decided, inside the text the agents reason over
(\S\ref{sec:leakage}). Detection is already implemented for the audit; removing
those sentences at extraction time would let the whole corpus support the claim
only the leak-free subset supports today, roughly doubling the evidence behind
this report's central accuracy figure at no research cost.

\textbf{Extend the remedy vocabulary beyond scoring challenges.} The corpus
contains three real remedies the system cannot express: a declaration of
ineffectiveness, financial penalties on the authority, and a court-determined
re-ranking. Each is a documented gap tied to a named case, not a hypothetical.

\textbf{Calibrate the risk screen against real outcomes.} This is the only thing
this project asks of Fusion21, and the highest-value item on the list. Anonymised
pre-action data from their member base (the profile fields the screen already
requests, plus whether a challenge followed) would let the severity-weighted
heuristic be replaced or validated by a fitted model, and the risk score reported as a calibrated
probability rather than an ordering. It is also the only component whose central
claim is directly falsifiable, which is why it deserves the data most.

The broader conclusion is narrower than the one I expected to write. A
multi-agent LLM system can conduct a procedurally coherent procurement dispute
negotiation, and can be constrained not to invent the arithmetic it reasons from.
What it cannot yet be shown to do is reason better than a single well-posed
prompt, and it can still be led into confident nonsense by a premise it was handed
rather than one it worked out. The value here sits less in the simulator than in
the evidence about what such a system does when nobody is checking, and in the
demonstration that finding that out requires screens built to catch a system that
is wrong in a fluent and plausible way.